% =============================================================================
% 研究档案 —— 《破衣烂衫者：患者1,260的日记》
% 批评性附录。非小说正文的一部分，不应与正文装订在一起。
% 用 xelatex 编译。
% =============================================================================
\documentclass[11pt,a4paper]{article}

\usepackage[margin=2.6cm]{geometry}
\usepackage{fontspec}
\usepackage{xeCJK}
\setmainfont{EB Garamond}
\setCJKmainfont{Noto Serif CJK SC}
\setCJKsansfont{Noto Serif CJK SC}
\usepackage{xcolor}
\usepackage{booktabs}
\usepackage{enumitem}
\usepackage{titlesec}
\usepackage[hidelinks]{hyperref}

\definecolor{sangue}{HTML}{8B0000}
\definecolor{tinta}{HTML}{2E241A}
\titleformat{\section}{\Large\bfseries\color{sangue}}{\thesection}{0.7em}{}
\titleformat{\subsection}{\large\bfseries\color{tinta}}{\thesubsection}{0.7em}{}
\setlist{itemsep=0.25em, topsep=0.35em}
\linespread{1.15}

\newcommand{\ressalva}[1]{%
  \par\vspace{0.5em}\noindent\fcolorbox{sangue!40}{sangue!4}{%
  \parbox{0.965\linewidth}{\small\textbf{方法说明。} #1}}\par\vspace{0.5em}}

\title{\vspace{-1.5cm}\textbf{《破衣烂衫者——患者1,260的日记》}\\[0.3em]
\large 研究档案：文学计划、技法、理论与限度}
\author{\small 三语版批评性附录}
\date{\small 2026年8月5日}

\begin{document}
\maketitle
\thispagestyle{empty}

\begin{abstract}
\noindent 本档案描述并分析《破衣烂衫者——患者1,260的日记》：一部以印刷超文本形式
构建的档案体小说，出版有三个单语版本和一个三语版本。档案汇集了实测的技术记录、
形式结构、所属的文学谱系、所运用的叙事技法、所调动的理论框架、历史依据，以及
包含已核实局限的批评分析。这是研究用附录，不是宣传材料：第\ref{sec:limites}节
与第\ref{sec:validacao}节的存在，正是为了陈述本作\textbf{未能}达成之处，以及
\textbf{尚未}获得外部验证之处。
\end{abstract}

\ressalva{本档案中的每一个数字均于2026年8月5日对作品源文件实测得出，而非估算。
凡属阐释性判断，文中明确标示。凡需要外部验证而该验证并不存在者——接受史、影响力、
比较性优劣——本档案选择不作断言，而非猜测。见第\ref{sec:validacao}节。}

% =============================================================================
\section{技术记录}

\begin{center}\small
\begin{tabular}{@{}ll@{}}
\toprule
\textbf{书名} & 破衣烂衫者——患者1,260的日记 \\
\textbf{作者} & Marcelo Dias de Carvalho Filho \\
\textbf{合著与插图} & Marcelo Claro \\
\textbf{ISBN} & 9798189170492 \\
\textbf{体裁} & 档案体小说；纪实性心理恐怖 \\
\textbf{形式} & 印刷超文本，84个碎片，216条阅读路线 \\
\textbf{语言} & 葡萄牙语（巴西）、英语、简体中文 \\
\textbf{印刷开本} & 160 $\times$ 230 毫米，平装 \\
\bottomrule
\end{tabular}
\end{center}

\vspace{0.6em}
\noindent\textbf{各版本实测篇幅}

\begin{center}\small
\begin{tabular}{@{}lrrr@{}}
\toprule
\textbf{版本} & \textbf{数字版} & \textbf{印刷版 160$\times$230} & \textbf{篇幅} \\
\midrule
葡萄牙语版 & 415 页 & 435 页 & 43,730 词 \\
英语版     & 411 页 & 427 页 & 44,978 词 \\
中文版     & 387 页 & 397 页 & 70,750 字 \\
三语版     & 1,061 页 & 1,115 页 & --- \\
\bottomrule
\end{tabular}
\end{center}

\noindent\small 中文与西方语言之间的词数不可直接比较；中文采用字数，这是通行的
计量方式。

% =============================================================================
\section{文学计划}

《破衣烂衫者》自我呈现为一份\textbf{法证档案}，而非一部小说。读者收到的是84个
碎片——回忆、鉴定报告、书信、录音转写、临床表格——由一位匿名档案管理员围绕一个
物件汇编而成：若阿金·安东尼奥·科雷亚的日记，在巴尔巴塞纳收容医院登记号为1,260。

从结构与副文本中可读出的计划包含三重运动：

\begin{enumerate}
\item \textbf{打捞一段被掩埋的历史。} 巴尔巴塞纳收容医院真实存在；史学估计，
二十世纪期间死于其高墙之内者数以万计，其中许多人根本没有任何精神科诊断。本作
建立在这片土壤之上。

\item \textbf{让形式承担论证。} 碎片化不是装饰：它演示了从文献残余中重建一个人
生命的不可能。读者必须自行拼合案卷，而拼合永远是局部的。

\item \textbf{将读者卷入。} 本作拒绝旁观者的位置。整整一族碎片（CONT）以第二人称
称呼读者，并将其登记入传递之链，编号为\textbf{患者1,263}。
\end{enumerate}

\subsection{传递之链}

核心装置是一列被占据身体的编号队列：

\begin{center}\small
\begin{tabular}{@{}cll@{}}
\toprule
\textbf{编号} & \textbf{何人} & \textbf{何处确立} \\
\midrule
1,259 & 若阿金之前的那一位 & CONT-07、CONT-08 \\
1,260 & 若阿金·安东尼奥·科雷亚 & 书名 \\
1,261 & 埃托尔·奥利维拉医生 & DOC-16（署名） \\
1,262 & 卢西亚·门德斯医生 & LUC-10、LUC-11、尾声 \\
\textbf{1,263} & \textbf{读者} & LUC-12、DOC-24、CONT-09、阅读协议 \\
\bottomrule
\end{tabular}
\end{center}

这一计数的精确性是装置生效的条件：正是它把最后一个位置从含糊的隐喻变成了针对
个人的威胁。

% =============================================================================
\section{形式结构}

\subsection{四个碎片族}

\begin{center}\small
\begin{tabular}{@{}llrrr@{}}
\toprule
\textbf{代号} & \textbf{性质} & \textbf{数量} & \textbf{词数} & \textbf{平均} \\
\midrule
MEM  & 患者1,260的回忆 & 27 & 13,310 & 492 \\
DOC  & 文件、鉴定报告、书信 & 27 & 15,247 & 564 \\
LUC  & 卢西亚·门德斯的调查（2026） & 17 & 8,934 & 525 \\
CONT & 读者污染协议 & 13 & 6,239 & 479 \\
\midrule
     & \textbf{合计} & \textbf{84} & \textbf{43,730} & 521 \\
\bottomrule
\end{tabular}
\end{center}

篇幅差异极大且具有功能性：从130词（MEM-26）到2,310词（DOC-03，尸检报告）。
短碎片如重击，长碎片则构成文献式沉浸。

\subsection{五部，三幕}

线性阅读穿越五个部分——塞尔唐（9）、收容医院（16）、奥利维拉的日记与鉴定报告
（25）、卢西亚的调查（20）、污染（14 + 尾声）——它们归并为三幕：创伤（9个碎片）、
体制化（41个）与循环（35个）。

这一不对称是有意义的：对应62年监禁的中间一幕，占据了全书近一半篇幅。监禁的
时长在阅读的时长中得到复现。

\subsection{路线网络}

\begin{center}\small
\begin{tabular}{@{}lr@{}}
\toprule
碎片 & 84 \\
阅读路线 & 216 \\
不同目的地 & 85（84个碎片 + 尾声） \\
无入线路线的碎片 & 0 \\
可抵达尾声的碎片 & 84 / 84 \\
循环（强连通分量 $>$ 1） & 6 \\
无出口的循环 & 0 \\
被引用最多的节点 & DOC-01，20条入线路线 \\
\bottomrule
\end{tabular}
\end{center}

两项属性是刻意为之并已核实：\textbf{一切路线通向尾声}，且\textbf{没有任何循环
困住读者}。循环被保留下来——它们正是本作的论题，文中称"循环重新开始"——但没有
一个是陷阱：从任何碎片出发，都存在通往终点的有限路径。

\ressalva{另一种方案是严格的无环图。该方案经测算后被否决：它将牺牲216条路线
中的85条，恰恰是那些赋予作品档案性质的"文件$\to$回忆"互指，并且会与文本自身的
论题相矛盾。此选择出自作者，并已记录在案。}

% =============================================================================
\section{文学与叙事谱系}

\ressalva{本节指出的是\textbf{可在文本中验证的形式亲缘性}。这并不断言作者读过
所引作品，也不断言其有意识地由此衍生。此处的"谱系"指手法上的亲缘——批评可加以
比较者——而非阅读的家谱。}

\subsection{档案体小说与伪造文献}

本作属于文学"伪纪录"传统，其现代母体是书信体：《德古拉》（1897）以日记、书信
与剪报拼装恐怖，并把任何角色都未能完成的综合工作交给读者。《破衣烂衫者》采用
同一原则，并通过纳入法证鉴定、保管链与保存状况评估——官僚文体而非私密文体——
将其激进化。

它也接近《叶之屋》（丹尼利维斯基，2000）：以书页为空间、设置导航装置、制造
书志学的眩晕。但有一处决定性差异：丹尼利维斯基把评注层级叠加成迷宫，
《破衣烂衫者》则让链条保持简短而具名，并用它抵达读者。

\subsection{印刷载体上的超文本小说}

《跳房子》（科塔萨尔，1963）确立了"一本书可有多种阅读顺序"的先例。
《破衣烂衫者》超出科塔萨尔的"棋盘"之处，在于提供了一个稠密网络（85个节点之间
216条边），其属性由自动化验证保证——收敛性与不困住读者。

\ressalva{该验证属于\textbf{制作}层面的特征，而非文学价值。它保证导航承诺得到
兑现；对散文品质则不置一词。}

\subsection{见证与暴行文学}

本作与巴西揭露收容所体制的传统对话，其当代新闻学标志是《巴西的大屠杀》
（达妮埃拉·阿尔贝克斯，2013），在DOC-19中被具名引用。在文学层面，第一部分
回响着《干枯的生命》（格拉西利亚诺·拉莫斯，1938）：干旱、举家迁徙、洗练的语言。

底层的理论问题是\textit{可再现性}：如何叙述暴行而不将其奇观化。本作的选择——
把超自然恐怖作为明确的隐喻，并在《历史注记》中如此声明——是一种立场，而且是
可争议的立场，第\ref{sec:critica}节对此加以检视。

\subsection{第二人称与被卷入的读者}

第二人称称呼在《如果在冬夜，一个旅人》（卡尔维诺，1979）中有经典先例，那里它
是游戏性与元文学的。《破衣烂衫者》反转了调性：不是与阅读嬉戏，而是指控阅读。
84个碎片中有27个持续动用"您"。

% =============================================================================
\section{叙事技法}

\subsection{诱导与暗示}

CONT系列碎片运用了临床暗示文献中可辨识的手法：

\begin{itemize}
\item \textbf{双重束缚。} CONT-03构造了一个无出路的选择："\textit{如果是：气味
就在那里。……如果不是：您在撒谎。您抬起来了。所有人都会抬。}" 两种回答都确证
前提。

\item \textbf{意动暗示。} "\textit{稍后，当您分神时，您的手指会自己回到书页上。}"

\item \textbf{后催眠暗示。} "\textit{今晚它会在您的枕头上。}"

\item \textbf{生理锚定。} 对症状的描述（呼吸、瞳孔、四肢温度）倾向于诱发那种
恰好证实症状的注意。

\item \textbf{定向空想性错视。} "\textit{看房间的右上角……盯着那个黑暗的角落
十秒钟。}"
\end{itemize}

\subsection{缺席的反转}

本作最精致的技法出现在CONT-04的"第四夜"：经历三夜噩梦之后，什么也没发生。而
文本把这份宽慰转化为威胁："\textit{恐惧的缺席比恐惧更糟……等待才是真正的传染。}"

此手法之所以值得注意，在于它撤除了刺激却保留了效果——把恐怖的生产从文本转移
到了读者身上。

\subsection{作为形式的临床量表}

DOC-09与CONT-05采用精神科自评量表的形式，设有条目、选项与最终计分。DOC-09的
第15项——"\textit{如果您现在能停止阅读并永不再打开这本书，您会停下吗？（）会
（）不会（）我已经试过}"——依靠双关运作：在小说内部指的是阅读；而在它所模仿的
临床工具语域中，读起来则是一个自杀风险条目。

\ressalva{这一歧义是有效的，也正因如此，卷首那页带有危机热线的警示并非装饰。
见第\ref{sec:etica}节。}

\subsection{边注与叙事内副文本}

本作使用从外部评注正文的边注、患者档案卡、档案封印，以及一整页完全属于叙事
内部的"当心"。副文本参与虚构，而非为其加框。

% =============================================================================
\section{所调动的理论框架}

\ressalva{本作未声明任何理论归属。以下是批评研究可以富有成效地加以运用的框架，
并指出作品中何者召唤了它们——而非把意图归给作者。}

\begin{description}[style=unboxed,leftmargin=0pt]
\item[创伤与代际传递。] 那个生物形象化了未经处理便穿越世代的痛苦。《历史注记》
本身即明言："\textit{那穿越世代、无人处理、如传染般传递下去的痛苦，才是唯一
真实的怪物}"。这是作品为自身提供的阅读钥匙。

\item[全控机构。] 收容医院的描写具备戈夫曼所指认的全控机构特征：生活各领域之间
屏障的瓦解、自我的贬损、特权的经济学。"被沉默者病区"是这一体制的强化。

\item[死亡政治与可弃置的生命。] 本作坚持指出，死于收容医院者是"穷人、黑人、
癫痫患者、单身母亲、同性恋者"。收容的标准是社会性的，而非临床性的——这把阅读
从精神医学的恐怖移向了政治的恐怖。

\item[档案、权力与空缺。] 带有文献空缺的保管链（DOC-20）与档案事故记录
（DOC-23）把档案主题化为权力机构，把文献的沉默主题化为附加的暴力。

\item[接受与被卷入的读者。] 本作构造了一个模范读者并为其命名。整个CONT族就是
一个被投入运作、而非被讨论的接受理论装置。
\end{description}

% =============================================================================
\section{历史依据与伦理立场}
\label{sec:etica}

本作的《历史注记》在方法上是审慎的，值得记录，因为这在同类作品中并不常见：

\begin{itemize}
\item 声明若阿金、那个生物以及被叙述的机构均为虚构。
\item 区分了虚构可能混淆的两个历史指涉：1915年的难民围栏（Alagadiço）与
塞纳多尔·蓬佩乌的围栏（1932—1933）。
\item 明确声明"\textit{收容人数与死亡总数并非以精确数据呈现}"。
\item 将巴尔巴塞纳收容医院的文献依据归于具名来源（阿尔贝克斯）。
\item 声明超自然恐怖是隐喻，而非论断。
\end{itemize}

书中还设有先于虚构的读者警示页，列举敏感内容、区分虚构与史实，并提供危机
热线。它属于副文本：位于虚构框架之前，不打断沉浸。

% =============================================================================
\section{批评分析}
\label{sec:critica}

\subsection{作品的强处}

\begin{enumerate}
\item \textbf{形式承担论证。} 碎片化、各部分之间的不对称以及路线网络并非装饰：
它们在形式上复现了重建的不可能与监禁的时长。

\item \textbf{读者的卷入是被执行的，而非被宣告的。} 编号之链终止于一个由读者
占据的空位。这是一个简单而构造严谨的装置。

\item \textbf{文献诚实是明示的。} 作品区分了何为虚构、何有文献依据，并拒绝
它并不拥有的精确性。

\item \textbf{诱导性散文在技术上是称职的。} CONT-04的反转是最佳例证：撤除刺激
而保留效果。
\end{enumerate}

\subsection{作品的可争议之处}

\begin{enumerate}
\item \textbf{超自然隐喻可能减轻历史恐怖。} 通过赋予苦难一个奇幻的原因，作品
有把政治责任转移到某个实体身上的风险。《历史注记》预见了这一反驳并作出回应；
该回应是否足够，是一个开放的批评问题，也是本作向讨论提供的最富成效的问题。

\item \textbf{编者注打断诱导。} 十三个碎片带有声明该装置属形式手法的脚注。其中
两处落在效果的顶点——CONT-03以"\textit{对您的身体而言，她是真实的}"收束，而
注释紧接着说明污染是文本性的。信息是必要的；位置是可争议的。将这些注释合并为
卷末的单一注记，可以同时保全两种功能。

\item \textbf{第二人称的密度可能造成饱和。} 27个碎片持续动用"您"。在连续的
线性阅读中，该手法趋于可预测——这一风险由路线阅读缓解，但印刷版本无法消除。
\end{enumerate}

% =============================================================================
\section{已核实的局限}
\label{sec:limites}

\begin{enumerate}
\item \textbf{三个版本并非同一本书。} 有十个碎片的结构差异超出语言差异所能解释
的范围。极端例子：CONT-05在葡语版与英语版为163行，中文版为63行；CONT-03在葡语版
为174行，英语版与中文版为102行；DOC-26在葡语版与英语版为110行，中文版为56行。
在三语版中，三个版本并排呈现，使差异对读者可见。\textbf{这是"国际参照性作品"
主张所面临的最严重局限。}

\item \textbf{封面美术低于印刷标准。} 封面与封底文件为1,046 $\times$ 1,504像素，
在印刷版面上仅得163 DPI；所需标准为300 DPI。

\item \textbf{封底带有测试用条形码。} 嵌入美术图中的元素标示ISBN
978-65-01-23456-7，并非本作的ISBN。

\item \textbf{语料未置于版本控制之下。} 作品源文件没有提交历史：既无差异比对，
也无法回退一次糟糕的编辑。
\end{enumerate}

% =============================================================================
\section{验证状态}
\label{sec:validacao}

本节的存在，是为了防止本档案被当作背书来读。

\begin{center}\small
\begin{tabular}{@{}p{6.2cm}l@{}}
\toprule
\textbf{维度} & \textbf{状态} \\
\midrule
五个版本的编译完整性 & 已核实 \\
路线的源文件与印刷件对等 & 已核实（648/648） \\
收敛性与不困住读者 & 已核实 \\
碎片完整性（SHA-256） & 已核实 \\
\midrule
学术同行评议 & \textbf{未进行} \\
外部专业编辑审读 & \textbf{未进行} \\
专家史学核查 & \textbf{未进行} \\
英译与中译的母语审读 & \textbf{未进行} \\
批评接受 & \textbf{无}（作品尚未出版） \\
影响力 & \textbf{无法测量} \\
\bottomrule
\end{tabular}
\end{center}

\ressalva{前四行属于\textbf{内部技术}验证：它们表明该物件构造良好，而非表明它
优秀。没有任何一项可以替代批评判断。后六项尚未进行，本档案也不预支它们。任何
由本文衍生的材料中出现的价值、原创性或影响力主张，在存在外部验证之前，均应
视为无根据。}

% =============================================================================
\section{研究议程}

本作开启、且研究可加以承接的问题：

\begin{enumerate}
\item 超自然隐喻究竟放大还是削弱了一部关于有据可查之暴行的叙事的见证力量？
\item 路线网络对阅读体验产生何种效应——以实证测量，而非以假设断言？
\item 贯穿27个碎片的第二人称称呼，在读者身上产生的是卷入还是抗拒？
\item 三个版本之间的差异如何改变各语言中被接受的作品——这又对"三语版"这一
概念本身说明了什么？
\item 保护性副文本（警示页、编者注）与作品所构建的沉浸装置如何相互作用？
\end{enumerate}

% =============================================================================
\section*{文中提及的作品}

\begin{itemize}\small
\item ARBEX, Daniela.《巴西的大屠杀》(Holocausto Brasileiro). Geração Editorial, 2013.
\item CALVINO, Italo.《如果在冬夜，一个旅人》. 1979.
\item CORTÁZAR, Julio.《跳房子》. 1963.
\item DANIELEWSKI, Mark Z.《叶之屋》(House of Leaves). Pantheon, 2000.
\item GOFFMAN, Erving.《收容所》(Asylums). 1961.
\item RAMOS, Graciliano.《干枯的生命》(Vidas Secas). 1938.
\item STOKER, Bram.《德古拉》(Dracula). 1897.
\end{itemize}

\vfill
\noindent\rule{\linewidth}{0.4pt}\\
{\footnotesize 本档案作为批评性附录生成。关于本项目度量与主张的保留说明，
参见 \texttt{CORRIGENDUM.md}。}

\end{document}
